\section{L'objectif du projet}
\subsection{La technologie $PUSH$}
Une nouvelle technologie appara\^{\i}t prometteuse pour Internet. Elle est
d�nomm�e {\it Push} et correspond \`a l'utilisation de son terminal
comme d'un r\'ecepteur, en continu, d'un flux d'information provenant de
diff\'erents serveurs. On ne va plus chercher l'information sur le serveur:
cette derni\`ere est directement d\'elivr\'ee sur la console du
r\'ecepteur. Cette technologie peut par exemple permettre de recevoir chez
soi les cours de la bourse, des d\'ep\^eches AFP, le bulletin
m\'et\'eo du lendemain, et surtout de suivre un cours ou une conf�rence.\\

Actuellement, l'on trouve sur Internet des techniques permettant ce type
d'applications, mais cela oblige les navigateurs \`a aller chercher, en
permanence, l'information sur les serveurs, et \`a la rapatrier. Il s'agit
donc en fait, \`a proprement parler, de {\it Pull} et non de {\it Push}.\\

De plus, avec les techniques existantes, le r�cepteur est connect� �
l'�metteur, ce qui du c�t� de ce dernier n�c�ssite la r�servation de
ressources, et limite donc le nombre de r�cepteurs.\\

L'on se propose, dans ce projet, d'utiliser le r\'eseau Multicast afin de
diffuser un flux d'information (en l'occurence, ici, une pr\'esentation de
transparents) d'un \'emetteur vers plusieurs r\'ecepteurs. Avec le
protocole que nous allons utiliser ({\it AMP: Asynchronous Multicast
Protocol}), le nombre de r�cepteurs est illimit� car il n'y a pas de
connexion �tablie avec l'�metteur). Et cette mani\`ere 
de proc\'eder correspond donc bien \`a la technologie {\it Push}.

\subsection{La diffusion de transparents}
L'\'emetteur poss\`ede le contr\^ole du flux d'information, c'est-\`a-dire
qu'il s'abonne \`a un canal Multicast. Il choisit les transparents et les
diffuse sur le canal IP Multicast dans l'ordre qui lui convient, avec le
rythme qu'il d\'esire (possibilit\'e d'aller en avant ou en arri\`ere, et
de rendre la s\'equence automatique).\\

Le r\'ecepteur, lui, ne peut que s'abonner \`a un canal Multicast, sur
lequel il recevra le flux d'information. Il ne peut pas exercer de
contr\^ole sur la r\'eception des transparents (si ce n'est d�marrer ou
interrompre la r�ception).\\

Les transparents se pr\'esentent sous la forme de pages HTML
classiques. Aussi, \'emetteur comme r\'ecepteur doivent disposer d'un
navigateur pour les visualiser. Le transfert sur le canal Multicast, par
contre, s'effectue sous la forme de blocs d'octets. Il y a donc conversion
des donn\'ees, entre leur utilisation et leur transport:
\begin{itemize}
  \item L'�metteur doit, avant �mission, convertir la page HTML en
    fichier binaire comprenant la page elle-m�me et les images, appelettes, et
    autres composants de la page.
  \item Le r�cepteur, lui, doit reconstituer, avant de l'afficher, la page
    HTML et les diff�rents fichiers auxquels fait r�f�rence la page, tout
    cela � partir du fichier binaire re\c cu.
\end{itemize}

Le projet a �t� r\'ealis\'e dans le langage de programmation Java.
Il se base sur le protocole de diffusion Multicast sur IP, {\it AMP
Asynchronous Multicast Protocol}, impl�ment� par St�phan Kraus � l'Institut
Eur�com.\\

De plus, il comporte plusieurs parties, correspondant aux trois couches
logiques suivantes:
\begin{htmllist}
  \htmlitemmark{RedBall}
  \item[Le contr\^ole:] d\'efilement de la s\'equence des
    transparents, et abonnement \`a Multicast,
  \item[La conversion des donn\'ees:] HTML $\rightarrow$
    Fichier Binaire, et Fichier Binaire $\rightarrow$ HTML, pour rendre
    les pages HTML transportables sur Multicast,
  \item[L'interface graphique:]
    \begin{htmllist}
      \htmlitemmark{WhiteBall}
      \item[Pour l'\'emetteur:] panel de contr\^ole de d\'efilement, gestion
        du nom de canal Multicast,
      \item[Pour le r\'ecepteur:] afficher les transparents, les enregistrer
        (fonction que nous n'avons pas impl�ment�e pour l'instant puisque les
        transparents sont visualis�s avec Netscape, et qu'il suffit donc
        d'effectuer un {\it Save As} pour enregister la page et ses composants
        depuis Netscape m�me) et les repasser ult\'erieurement (toujours en
        utilisant Netscape ou un autre �diteur de pages HTML).
    \end{htmllist}
\end{htmllist}
Bien s�r, il y a aussi une quatri�me couche, celle de transport des fichiers
binaires: c'est celle qui impl�mente AMP.
